\documentclass[11pt,a4paper]{article}

\usepackage[margin=25mm]{geometry}
\usepackage[T1]{fontenc}
\usepackage{lmodern}
\usepackage{amsmath}
\usepackage{graphicx}
\usepackage{siunitx}
\usepackage{booktabs}
\usepackage{float}
\usepackage[hidelinks]{hyperref}

\newcommand{\manual}[1]{(Laboratory manual, p.~#1)}

\begin{document}


\title{Laboratory Exercise 1: Dynamical Pendulum}
\author{Magnus Støren Wedèn, Jon Magnus Vårheim, Birk Holm Broen}
\date{September 2026}

\maketitle

\section{Introduction}
Laboratory exercise 1 will analyse a cart-pendulum system to study equilibrium, stability and feedback control. There are three sections to this lab, stable point 1 where the controller is not active, this causes the pendulum to hang downwards due to gravity, stable point 2 occurs when the controller is active attempting to force the tip of the pendulum stay upright and finally the situation where the controller is active and exposed to outside disturbance. 


\section{Method}

\section{First watch observations}
An initial qualitative review of the video was used to identify
the main phases of the experiment before taking detailed
measurements. The review focused on the overall motion of the
pendulum and cart, changes in operating condition, and possible
disturbance events. This follows the observation procedure
described in the laboratory manual. \manual{3--4}

The recording was divided into intervals covering natural motion
around Stable Point 1 (SP1), controlled operation around Stable
Point 2 (SP2), and disturbances during upright operation.
SP1 denotes the naturally stable downward equilibrium, whereas
SP2 denotes the upright equilibrium maintained by feedback
control. \manual{2, 6--8}

For the SP1 intervals, the analysis considers displacement and
release of the pendulum, its subsequent oscillation about the
downward equilibrium, and any visible cart motion.
For the SP2 intervals, the analysis considers angular deviations
from upright and the associated corrective cart movements.
The disturbance interval is examined more closely to identify
changes in motion, the sequence of visible responses, and the
subsequent recovery or loss of stabilisation.
\manual{10--11, 14--20}

Table~\ref{tab:intervals} lists the selected analysis intervals.
Timestamps refer to the original recording and use the format
hh:mm:ss.sss. The timestamps and frame boundaries were obtained
from the segment exports and matched by segment name.
The interval descriptions identify the intended focus of the
analysis; they do not establish the cause of an observed event.


\begin{table}[H]
    \centering
    \small
    \caption{Selected analysis intervals in the original video,
    with timestamps and exported frame boundaries.}
    
    \label{tab:intervals}
    \begin{tabular}{@{}lccrr@{}}
        \toprule
        Analysis interval
        & Start time
        & End time
        & \shortstack{Start\\frame}
        & \shortstack{End\\frame} \\
        \midrule
        SP1: initial observation
        & 00:00:08.041 & 00:00:11.812 & 241 & 354 \\

        SP1: input/disturbance
        & 00:00:12.980 & 00:00:15.452 & 389 & 463 \\

        SP1: input/disturbance continued
        & 00:00:20.287 & 00:00:23.790 & 608 & 713 \\

        SP1: oscillation
        & 00:00:33.188 & 00:00:38.885 & 995 & 1165 \\

        SP2: controlled operation 1
        & 00:00:53.020 & 00:01:09.937 & 1589 & 2096 \\

        SP2: controlled operation 2
        & 00:01:13.974 & 00:01:22.616 & 2217 & 2476 \\

        SP2: disturbance sequence
        & 00:01:54.015 & 00:02:01.197 & 3417 & 3632 \\
        \bottomrule
    \end{tabular}
\end{table}

Record video frame rate $f_\text{video}$ and frame duration $\Delta t_\text{frame} = 1/f_\text{video}$\newline

\begin{gather}
t = 00{:}00{:}08.041 = 8.041~\text{s} \nonumber \\
f_\text{video} = \frac{N}{t} \label{eq:fvideo} \\
f_\text{video} = \frac{241}{8.041~\text{s}} \approx 29.97~\text{frames/s} \nonumber \\
\Delta t_\text{frame} = \frac{1}{f_\text{video}} \label{eq:dtframe} \\
\Delta t_\text{frame} = \frac{1}{29.97~\text{frames/s}} \approx 0.0334~\text{s} \nonumber
\end{gather}


\subsection{Candidate disturbance-event windows}
\label{sec:disturbance-windows}

Within the selected disturbance sequence, five candidate
event windows were marked for closer inspection, as listed
in Table~\ref{tab:disturbance-windows}.
These windows guide the frame-by-frame analysis.
They are preliminary event locations, rather than confirmed
disturbance durations or onset-time uncertainty intervals.

For each candidate event, direct observations are recorded
before a disturbance mechanism is proposed. The analysis
considers which component responds first, whether corrective
cart motion remains visible, and whether the pendulum returns
to upright operation. Where the evidence does not support a
reliable explanation, the disturbance is classified as unknown.
\manual{3--4, 17--20}

\begin{table}[htbp]
    \centering
    \small
    \caption{Preliminary disturbance-event windows in the
    original video. Timestamps use hh:mm:ss.sss.}
    \label{tab:disturbance-windows}
    \begin{tabular}{@{}ccc@{}}
        \toprule
        Candidate event & Window start & Window end \\
        \midrule
        D1 & 00:01:55.215 & 00:01:55.248 \\
        D2 & 00:01:56.083 & 00:01:56.116 \\
        D3 & 00:01:56.850 & 00:01:56.950 \\
        D4 & 00:01:57.784 & 00:01:57.884 \\
        D5 & 00:02:00.354 & 00:02:00.454 \\
        \bottomrule
    \end{tabular}
\end{table}

\section{Settle point 1}
\begin{table}[H]
    \centering
    \caption{SP1: all four measurements}
    \label{tab:sp1_all}
    \renewcommand{\arraystretch}{1.6}
    \begin{tabular}{|l|c|c|c|c|c|}
        \hline
        \textbf{Quantity} & \textbf{Symbol} & \multicolumn{4}{c|}{\textbf{Measurement}} \\
        \cline{3-6}
         & & \textbf{1} & \textbf{2} & \textbf{3} & \textbf{4} \\
        \hline
        Release time (s)                        & $t_0$              & 33.000 & 32.933 & 32.799 & 33.066 \\ \hline
        Initial angular displacement ($^\circ$) & $\tilde{\theta}_0$ & 170.1  & 171.4  & 172.0  & 169.2  \\ \hline
        First selected peak ($^\circ$)          & $\tilde{\theta}_1$ & 141.7  & 144.7  & 141.4  & 142.3  \\ \hline
        Second selected peak ($^\circ$)         & $\tilde{\theta}_2$ & 30.2   & 30.0   & 27.6   & 28.2   \\ \hline
        Third selected peak ($^\circ$)          & $\tilde{\theta}_3$ & 157.2  & 158.6  & 157.6  & 159.8  \\ \hline
        Time of first selected peak (s)         & $t_1$              & 34.401 & 34.434 & 34.468 & 34.368 \\ \hline
        Time of second selected peak (s)        & $t_2$              & 35.235 & 35.202 & 35.269 & 35.168 \\ \hline
        Time of third selected peak (s)         & $t_3$              & 35.936 & 36.903 & 35.969 & 36.003 \\ \hline
        Settling time (s)                       & $t_s$              & not visible & not visible & not visible  &    not visible    \\ \hline
    \end{tabular}
\end{table}


\begin{align}
T_i &= t_{i+1} - t_i \label{eq:Ti} \\
\bar{T} &= \frac{1}{n}\sum_{i=1}^{n} T_i \label{eq:Tmean} \\
f_d &= \frac{1}{\bar{T}} \label{eq:fd} \\
\omega_d &= 2\pi f_d \label{eq:wd} \\
r_i &= \frac{|\tilde{\theta}_{i+1}|}{|\tilde{\theta}_i|} \label{eq:ri} \\
\bar{r} &= \frac{1}{n_r}\sum_{i=1}^{n_r} r_i \label{eq:rmean} \\
\bar{x} &= \frac{\sum x_i}{N} \label{eq:mean} \\
SD_\text{sample} &= \sqrt{\frac{\sum \left(x_i-\bar{x}\right)^2}{N-1}} \label{eq:sd}
\end{align}



Example calculation for measurement 1:
\begin{align*}
T_1 &= 35.235 - 34.401 = 0.834~\text{s} \\
T_2 &= 35.936 - 35.235 = 0.701~\text{s} \\
\bar{T} &= \frac{0.834 + 0.701}{2} = 0.768~\text{s} \\
f_d &= \frac{1}{0.768~\text{s}} = 1.303~\text{Hz} \\
\omega_d &= 2\pi \cdot 1.303~\text{Hz} = 8.187~\text{rad/s} \\
r_1 &= \frac{30.2}{141.7} = 0.213 \\
r_2 &= \frac{157.2}{30.2} = 5.205 \\
\bar{r} &= \frac{0.213 + 5.205}{2} = 2.709
\end{align*}

Example standard deviation, mean period over the four measurements:
\begin{align*}
\bar{T} &= \frac{0.768 + 1.235 + 0.751 + 0.818}{4} = 0.893~\text{s} \\
SD_{T} &= \sqrt{\frac{(0.768-0.893)^2 + (1.235-0.893)^2 + (0.751-0.893)^2 + (0.818-0.893)^2}{4-1}} \\
SD_{T} &= \sqrt{\frac{0.1584}{3}} = 0.230~\text{s}
\end{align*}

\begin{table}[H]
    \centering
    \caption{SP1: final values, given as mean $\pm$ sample standard deviation ($N = 4$).}
    \label{tab:sp1_results}
    \renewcommand{\arraystretch}{1.6}
    \begin{tabular}{|l|c|c|}
        \hline
        \textbf{Quantity} & \textbf{Symbol} & \textbf{Value} \\
        \hline
        Mean period                    & $\bar{T}$  & $0.893 \pm 0.230$~s \\ \hline
        Observed oscillation frequency & $f_d$      & $1.167 \pm 0.243$~Hz \\ \hline
        Damped angular frequency       & $\omega_d$ & $7.334 \pm 1.524$~rad/s \\ \hline
        Mean amplitude ratio           & $\bar{r}$  & $2.835 \pm 0.125$ \\ \hline
    \end{tabular}
\end{table}

\section{Settle point 2}
\begin{table}[H]
    \centering
    \caption{SP2: all four measurements}
    \label{tab:sp2_all}
    \renewcommand{\arraystretch}{1.6}
    \begin{tabular}{|l|c|c|c|c|c|c|c|}
        \hline
        \textbf{Quantity} & \textbf{Symbol} & \multicolumn{4}{c|}{\textbf{Measurement}} & \textbf{Mean} & \textbf{s} \\
        \cline{3-6}
         & & \textbf{1} & \textbf{2} & \textbf{3} & \textbf{4} & & \\
        \hline
        Beginning of controlled operation (s) & $t_\text{on}$               & 53.020 & 53.053 & 52.986 & 52.953 & 53.003 & 0.043 \\ \hline
        Maximum positive deviation ($^\circ$) & $\tilde{\theta}_\text{max}$ & 4.1    & 2.3    & 3.2    & 4.9    & 3.6    & 1.1   \\ \hline
        Maximum negative deviation ($^\circ$) & $\tilde{\theta}_\text{min}$ & $-0.9$ & $-3.8$ & $-2.8$ & $-2.6$ & $-2.5$ & 1.2   \\ \hline
        Maximum cart position (px)            & $x_\text{max}$              & 1184   & 1181   & 1186   & 1183   & 1183.5 & 2.1   \\ \hline
        Minimum cart position (px)            & $x_\text{min}$              & 765.8  & 759.0  & 749.5  & 755.5  & 757.5  & 6.8   \\ \hline
        Stabilisation time (s)                & $t_\text{stab}$             & 0.334  & 2.369  & 3.170  & 1.501  & 1.84   & 1.22  \\ \hline
    \end{tabular}
\end{table}

% ---------- Results section ----------
Example mean and standard deviation, beginning of controlled operation, using Eqs.~\eqref{eq:mean} and \eqref{eq:sd}:
\begin{align*}
\bar{t}_\text{on} &= \frac{53.020 + 53.053 + 52.986 + 52.953}{4} = 53.003~\text{s} \\
SD_{t_\text{on}} &= \sqrt{\frac{(53.020-53.003)^2 + (53.053-53.003)^2 + (52.986-53.003)^2 + (52.953-53.003)^2}{4-1}} \\
SD_{t_\text{on}} &= \sqrt{\frac{0.005578}{3}} = 0.043~\text{s}
\end{align*}

% ---------- Results section ----------
\begin{align}
\theta_\text{range} &= \tilde{\theta}_\text{max} - \tilde{\theta}_\text{min} \label{eq:thrange} \\
\tilde{\theta}_\text{abs,max} &= \max_t |\tilde{\theta}(t)| \label{eq:thabs} \\
x_\text{range} &= x_\text{max} - x_\text{min} \label{eq:xrange}
\end{align}

Example calculation for measurement 1:
\begin{align*}
\theta_\text{range} &= 4.1^\circ - (-0.9^\circ) = 5.0^\circ \\
\tilde{\theta}_\text{abs,max} &= \max\left(|4.1^\circ|, |-0.9^\circ|\right) = 4.1^\circ \\
x_\text{range} &= 1184~\text{px} - 765.8~\text{px} = 418.2~\text{px}
\end{align*}

Example standard deviation, angular operating range over the four measurements, using Eq.~\eqref{eq:sd}:
\begin{align*}
\bar{\theta}_\text{range} &= \frac{5.0 + 6.1 + 6.0 + 7.5}{4} = 6.15^\circ \\
SD_{\theta_\text{range}} &= \sqrt{\frac{(5.0-6.15)^2 + (6.1-6.15)^2 + (6.0-6.15)^2 + (7.5-6.15)^2}{4-1}} \\
SD_{\theta_\text{range}} &= \sqrt{\frac{3.17}{3}} = 1.03^\circ
\end{align*}

\begin{table}[H]
    \centering
    \caption{SP2: final values, given as mean $\pm$ sample standard deviation ($N = 4$).}
    \label{tab:sp2_results}
    \renewcommand{\arraystretch}{1.6}
    \begin{tabular}{|l|c|c|}
        \hline
        \textbf{Quantity} & \textbf{Symbol} & \textbf{Value} \\
        \hline
        Angular operating range        & $\theta_\text{range}$          & $6.15 \pm 1.03^\circ$ \\ \hline
        Maximum absolute deviation     & $\tilde{\theta}_\text{abs,max}$ & $4.00 \pm 0.71^\circ$ \\ \hline
        Cart-position range            & $x_\text{range}$               & $426.1 \pm 7.9$~px \\ \hline
    \end{tabular}
\end{table}

\subsection{Linearisation around SP2}

The simplified nonlinear equation for the pendulum near the upright position is
\begin{equation}
J_p\ddot{\theta} + b_p\dot{\theta} - mgl\sin\theta = -ml\ddot{x}\cos\theta.
\label{eq:nonlinear}
\end{equation}
The constants are $m = 0.120$~kg, $J_p = 0.0139231$~kg\,m$^2$,
$b_p = 1.07443\times10^{-4}$~N\,m\,s/rad and $g = 9.81$~m/s$^2$
\manual{9}. The distance $l$ is not given. Since
$J_p = J_\text{cm} + ml^2 \ge ml^2$, it is bounded by
\begin{equation*}
l \le \sqrt{\frac{J_p}{m}} = \sqrt{\frac{0.0139231}{0.120}} = 0.341~\text{m},
\end{equation*}
and this upper value is used as an estimate below.

\paragraph{Step 1: Taylor expansions.}
First-order Taylor expansions around $\hat{\theta} = 0$ give
\begin{align}
\sin\theta &\approx \sin\hat{\theta} + \cos\hat{\theta}\,(\theta - \hat{\theta}) = \tilde{\theta}, \label{eq:sin_lin} \\
\cos\theta &\approx \cos\hat{\theta} - \sin\hat{\theta}\,(\theta - \hat{\theta}) = 1. \label{eq:cos_lin}
\end{align}
At the largest measured deviation, $4.9^\circ = 0.0855$~rad (Table~\ref{tab:sp2_all}),
\begin{align*}
\frac{\tilde{\theta} - \sin\tilde{\theta}}{\sin\tilde{\theta}} &= \frac{0.0855 - 0.0854}{0.0854} \approx 0.12\,\%, \\
\frac{1 - \cos\tilde{\theta}}{\cos\tilde{\theta}} &= \frac{1 - 0.9963}{0.9963} \approx 0.37\,\%.
\end{align*}

\paragraph{Step 2: Substitution.}
With $\theta = \hat{\theta} + \tilde{\theta}$ and $x = \hat{x} + \tilde{x}$, the derivatives become
\begin{equation*}
\dot{\theta} = \dot{\hat{\theta}} + \dot{\tilde{\theta}}, \qquad
\ddot{\theta} = \ddot{\hat{\theta}} + \ddot{\tilde{\theta}}, \qquad
\ddot{x} = \ddot{\hat{x}} + \ddot{\tilde{x}}.
\end{equation*}
Inserting these and Eqs.~\eqref{eq:sin_lin}--\eqref{eq:cos_lin} into Eq.~\eqref{eq:nonlinear} gives
\begin{equation*}
J_p\left(\ddot{\hat{\theta}} + \ddot{\tilde{\theta}}\right)
+ b_p\left(\dot{\hat{\theta}} + \dot{\tilde{\theta}}\right)
- mgl\left(\sin\hat{\theta} + \cos\hat{\theta}\,\tilde{\theta}\right)
= -ml\left(\ddot{\hat{x}} + \ddot{\tilde{x}}\right)\left(\cos\hat{\theta} - \sin\hat{\theta}\,\tilde{\theta}\right).
\end{equation*}

\paragraph{Step 3: Equilibrium terms.}
At the operating point, $\hat{\theta} = \dot{\hat{\theta}} = \ddot{\hat{\theta}} = 0$, and $\hat{x} = x_\text{ref}$ is constant, so $\ddot{\hat{x}} = 0$.
The equilibrium equation reduces to $0 = 0$, and all operating-point terms drop out:
\begin{equation*}
J_p\ddot{\tilde{\theta}} + b_p\dot{\tilde{\theta}} - mgl\,\tilde{\theta}
= -ml\,\ddot{\tilde{x}}\left(1 - 0\cdot\tilde{\theta}\right).
\end{equation*}

\paragraph{Step 4: Small deviations.}
The remaining product $\ddot{\tilde{x}}\,\tilde{\theta}$ is of second order in the small deviations and is multiplied by $\sin\hat{\theta} = 0$, so it is neglected.

\paragraph{Step 5: Linearised equation.}
The linearised small-deviation equation is
\begin{equation}
J_p\ddot{\tilde{\theta}} + b_p\dot{\tilde{\theta}} - mgl\,\tilde{\theta} = -ml\,\ddot{\tilde{x}}.
\label{eq:linearised}
\end{equation}
With $mgl = 0.120 \cdot 9.81 \cdot 0.341 = 0.401$~N\,m and $ml = 0.120 \cdot 0.341 = 0.0409$~kg\,m,
\begin{equation*}
0.01392\,\ddot{\tilde{\theta}} + 1.074\times10^{-4}\,\dot{\tilde{\theta}} - 0.401\,\tilde{\theta} = -0.0409\,\ddot{\tilde{x}},
\end{equation*}
or, solved for the angular acceleration,
\begin{equation*}
\ddot{\tilde{\theta}} = 28.8\,\tilde{\theta} - 0.0077\,\dot{\tilde{\theta}} - 2.94\,\ddot{\tilde{x}}.
\end{equation*}

\paragraph{Step 6: Instability without control.}
Without control, $\ddot{\tilde{x}} = 0$, and the gravitational term $+28.8\,\tilde{\theta}$ has the same sign as $\tilde{\theta}$.
Any deviation therefore produces an angular acceleration away from the upright position.
The characteristic equation of Eq.~\eqref{eq:linearised} with $\ddot{\tilde{x}} = 0$ is
\begin{equation}
J_p s^2 + b_p s - mgl = 0.
\label{eq:char}
\end{equation}
With the values above,
\begin{align*}
s &= \frac{-b_p \pm \sqrt{b_p^2 + 4J_p\,mgl}}{2J_p}
= \frac{-1.074\times10^{-4} \pm \sqrt{0.02233}}{2 \cdot 0.01392}, \\
s_1 &= 5.36~\text{s}^{-1}, \qquad s_2 = -5.37~\text{s}^{-1}.
\end{align*}
The positive root $s_1$ means the deviation grows as $e^{5.36t}$ and doubles every
\begin{equation*}
t_2 = \frac{\ln 2}{s_1} = \frac{0.693}{5.36~\text{s}^{-1}} = 0.13~\text{s}.
\end{equation*}

\paragraph{Step 7: Stabilisation by cart acceleration.}
The term $-\left(ml/J_p\right)\ddot{\tilde{x}}$ allows the cart to counteract gravity.
If the cart accelerates in proportion to the angular deviation, $\ddot{\tilde{x}} = k\,\tilde{\theta}$, Eq.~\eqref{eq:linearised} becomes
\begin{equation}
J_p\ddot{\tilde{\theta}} + b_p\dot{\tilde{\theta}} + \left(mlk - mgl\right)\tilde{\theta} = 0.
\label{eq:controlled}
\end{equation}
For $k > g = 9.81$~m/s$^2$ per radian, the coefficient of $\tilde{\theta}$ becomes positive and the net effect is restoring.
At the mean maximum deviation $\tilde{\theta}_\text{abs,max} = 4.00^\circ = 0.0698$~rad (Table~\ref{tab:sp2_results}), this requires a cart acceleration of at least
\begin{equation*}
\ddot{\tilde{x}} > g\,\tilde{\theta} = 9.81 \cdot 0.0698 = 0.68~\text{m/s}^2.
\end{equation*}
The cart thus moves under the pendulum in the direction in which it is falling.

\section{SP2 with disturbances: }
\begin{table}[H]
    \centering
    \caption{Disturbance SP2: all four measurements. No.\ = measurement.event.
    $u_\theta$ and $u_x$ are the sample standard deviations of the four
    measurements of each event.}
    \label{tab:sp2_disturbance}
    \renewcommand{\arraystretch}{1.6}
    \begin{tabular}{|c|c|c|c|c|}
        \hline
        No.
        & $t_d \pm u_{t_d}$ (s)
        & $\tilde{\theta}_\text{max} \pm u_\theta$ ($^\circ$)
        & $\Delta x_\text{max} \pm u_x$ (px)
        & $t_r \pm u_{t_r}$ (s) \\
        \hline
        1.1 & $115.182 \pm 0.033$ & $3.2 \pm 0.60$ & $75.3 \pm 11.6$  & not achieved \\ \hline
        2.1 & $115.182 \pm 0.033$ & $3.5 \pm 0.60$ & $94.2 \pm 11.6$  & not achieved \\ \hline
        3.1 & $115.215 \pm 0.033$ & $3.3 \pm 0.60$ & $89.8 \pm 11.6$  & not achieved \\ \hline
        4.1 & $115.182 \pm 0.033$ & $4.5 \pm 0.60$ & $103.1 \pm 11.6$ & not achieved \\ \hline\hline
        1.2 & $116.049 \pm 0.033$ & $6.7 \pm 1.28$ & $92.9 \pm 5.0$   & not achieved \\ \hline
        2.2 & $116.016 \pm 0.033$ & $4.1 \pm 1.28$ & $89.3 \pm 5.0$   & not achieved \\ \hline
        3.2 & $116.049 \pm 0.033$ & $6.4 \pm 1.28$ & $99.1 \pm 5.0$   & not achieved \\ \hline
        4.2 & $116.049 \pm 0.033$ & $6.8 \pm 1.28$ & $99.6 \pm 5.0$   & not achieved \\ \hline\hline
        1.3 & $116.850 \pm 0.033$ & $5.9 \pm 0.43$ & $104.3 \pm 3.2$  & not achieved \\ \hline
        2.3 & $116.917 \pm 0.033$ & $6.2 \pm 0.43$ & $98.2 \pm 3.2$   & not achieved \\ \hline
        3.3 & $116.883 \pm 0.033$ & $6.5 \pm 0.43$ & $99.9 \pm 3.2$   & not achieved \\ \hline
        4.3 & $117.883 \pm 0.033$ & $5.5 \pm 0.43$ & $104.5 \pm 3.2$  & not achieved \\ \hline\hline
        1.4 & $117.784 \pm 0.033$ & $5.5 \pm 0.55$ & $59.6 \pm 16.5$  & $1.001 \pm 0.047$ \\ \hline
        2.4 & $117.784 \pm 0.033$ & $4.4 \pm 0.55$ & $60.9 \pm 16.5$  & $1.035 \pm 0.047$ \\ \hline
        3.4 & $117.851 \pm 0.033$ & $4.9 \pm 0.55$ & $94.4 \pm 16.5$  & $0.934 \pm 0.047$ \\ \hline
        4.4 & $117.784 \pm 0.033$ & $4.3 \pm 0.55$ & $64.1 \pm 16.5$  & $1.001 \pm 0.047$ \\ \hline
    \end{tabular}
\end{table}

\begin{align}
\tilde{\theta}_\text{max} &= \max_{t \ge t_d} |\tilde{\theta}(t)| \label{eq:dthmax} \\
\Delta x_\text{max} &= \max_{t \ge t_d} |x(t) - \hat{x}_0| \label{eq:dxmax} \\
t_r &= t_\text{enter} - t_d \label{eq:tr} \\
u_{t_r} &= \sqrt{u_{t_\text{enter}}^2 + u_{t_d}^2} \label{eq:utr}
\end{align}

Example calculation for measurement 1, event 4:
\begin{align*}
t_r &= 118.785~\text{s} - 117.784~\text{s} = 1.001~\text{s} \\
u_{t_r} &= \sqrt{(0.033~\text{s})^2 + (0.033~\text{s})^2} = 0.047~\text{s}
\end{align*}

Example standard deviation, maximum angular deviation of event 1, using Eq.~\eqref{eq:sd}:
\begin{align*}
\bar{\tilde{\theta}}_{\text{max},1} &= \frac{3.2 + 3.5 + 3.3 + 4.5}{4} = 3.63^\circ \\
SD_{\theta,1} &= \sqrt{\frac{(3.2-3.63)^2 + (3.5-3.63)^2 + (3.3-3.63)^2 + (4.5-3.63)^2}{4-1}} \\
SD_{\theta,1} &= \sqrt{\frac{1.0675}{3}} = 0.60^\circ
\end{align*}

\begin{table}[H]
    \centering
    \caption{Disturbance SP2: final values per event, given as mean $\pm$ sample standard deviation ($N = 4$).}
    \label{tab:dist_results}
    \renewcommand{\arraystretch}{1.6}
    \begin{tabular}{|c|c|c|c|c|}
        \hline
        \textbf{Event} & $t_d$ (s) & $\tilde{\theta}_\text{max}$ ($^\circ$) & $\Delta x_\text{max}$ (px) & $t_r$ (s) \\
        \hline
        1 & $115.190 \pm 0.017$ & $3.63 \pm 0.60$ & $90.6 \pm 11.6$  & not achieved \\ \hline
        2 & $116.041 \pm 0.017$ & $6.00 \pm 1.28$ & $95.2 \pm 5.0$   & not achieved \\ \hline
        3 & $117.133 \pm 0.501$ & $6.03 \pm 0.43$ & $101.7 \pm 3.2$  & not achieved \\ \hline
        4 & $117.801 \pm 0.034$ & $4.78 \pm 0.55$ & $69.8 \pm 16.5$  & $0.993 \pm 0.042$ \\ \hline
    \end{tabular}
\end{table}



\section{Appendix}
\noindent{\large\bfseries Calculations for stable point 1:}\par
\begin{align*}
&\textbf{1:} & T_1 &= 35.235-34.401 = 0.834~\text{s}, & T_2 &= 35.936-35.235 = 0.701~\text{s} \\
& & \bar{T} &= \frac{0.834+0.701}{2} = 0.768~\text{s}, & f_d &= \frac{1}{0.768} = 1.303~\text{Hz} \\
& & \omega_d &= 2\pi \cdot 1.303 = 8.187~\text{rad/s}, & & \\
& & r_1 &= \frac{30.2}{141.7} = 0.213, & r_2 &= \frac{157.2}{30.2} = 5.205 \\
& & \bar{r} &= \frac{0.213+5.205}{2} = 2.709 & & \\[6pt]
&\textbf{2:} & T_1 &= 35.202-34.434 = 0.768~\text{s}, & T_2 &= 36.903-35.202 = 1.701~\text{s} \\
& & \bar{T} &= \frac{0.768+1.701}{2} = 1.235~\text{s}, & f_d &= \frac{1}{1.235} = 0.810~\text{Hz} \\
& & \omega_d &= 2\pi \cdot 0.810 = 5.090~\text{rad/s}, & & \\
& & r_1 &= \frac{30.0}{144.7} = 0.207, & r_2 &= \frac{158.6}{30.0} = 5.287 \\
& & \bar{r} &= \frac{0.207+5.287}{2} = 2.747 & & \\[6pt]
&\textbf{3:} & T_1 &= 35.269-34.468 = 0.801~\text{s}, & T_2 &= 35.969-35.269 = 0.700~\text{s} \\
& & \bar{T} &= \frac{0.801+0.700}{2} = 0.751~\text{s}, & f_d &= \frac{1}{0.751} = 1.332~\text{Hz} \\
& & \omega_d &= 2\pi \cdot 1.332 = 8.372~\text{rad/s}, & & \\
& & r_1 &= \frac{27.6}{141.4} = 0.195, & r_2 &= \frac{157.6}{27.6} = 5.710 \\
& & \bar{r} &= \frac{0.195+5.710}{2} = 2.953 & & \\[6pt]
&\textbf{4:} & T_1 &= 35.168-34.368 = 0.800~\text{s}, & T_2 &= 36.003-35.168 = 0.835~\text{s} \\
& & \bar{T} &= \frac{0.800+0.835}{2} = 0.818~\text{s}, & f_d &= \frac{1}{0.818} = 1.223~\text{Hz} \\
& & \omega_d &= 2\pi \cdot 1.223 = 7.686~\text{rad/s}, & & \\
& & r_1 &= \frac{28.2}{142.3} = 0.198, & r_2 &= \frac{159.8}{28.2} = 5.667 \\
& & \bar{r} &= \frac{0.198+5.667}{2} = 2.932 & &
\end{align*}

\begin{align*}
\bar{T} &= \frac{0.768+1.235+0.751+0.818}{4} = 0.893~\text{s}, &
SD_{T} &= \sqrt{\frac{0.1584}{4-1}} = 0.230~\text{s} \\
\bar{f}_d &= \frac{1.303+0.810+1.332+1.223}{4} = 1.167~\text{Hz}, &
SD_{f_d} &= \sqrt{\frac{0.1764}{4-1}} = 0.243~\text{Hz} \\
\bar{\omega}_d &= \frac{8.187+5.090+8.372+7.686}{4} = 7.334~\text{rad/s}, &
SD_{\omega_d} &= \sqrt{\frac{6.965}{4-1}} = 1.524~\text{rad/s} \\
\bar{r} &= \frac{2.709+2.747+2.953+2.932}{4} = 2.835, &
SD_{r} &= \sqrt{\frac{0.0469}{4-1}} = 0.125
\end{align*}

% ---------- Appendix ----------
\noindent{\large\bfseries SP2 calculations:}\par

\begin{align*}
&\textbf{1:} & \theta_\text{range} &= 4.1-(-0.9) = 5.0^\circ, & \tilde{\theta}_\text{abs,max} &= 4.1^\circ, & x_\text{range} &= 1184-765.8 = 418.2~\text{px} \\
&\textbf{2:} & \theta_\text{range} &= 2.3-(-3.8) = 6.1^\circ, & \tilde{\theta}_\text{abs,max} &= 3.8^\circ, & x_\text{range} &= 1181-759.0 = 422.0~\text{px} \\
&\textbf{3:} & \theta_\text{range} &= 3.2-(-2.8) = 6.0^\circ, & \tilde{\theta}_\text{abs,max} &= 3.2^\circ, & x_\text{range} &= 1186-749.5 = 436.5~\text{px} \\
&\textbf{4:} & \theta_\text{range} &= 4.9-(-2.6) = 7.5^\circ, & \tilde{\theta}_\text{abs,max} &= 4.9^\circ, & x_\text{range} &= 1183-755.5 = 427.5~\text{px}
\end{align*}

\begin{align*}
\bar{\theta}_\text{range} &= \frac{5.0+6.1+6.0+7.5}{4} = 6.15^\circ, &
SD_{\theta_\text{range}} &= \sqrt{\frac{3.17}{4-1}} = 1.03^\circ \\
\bar{\tilde{\theta}}_\text{abs,max} &= \frac{4.1+3.8+3.2+4.9}{4} = 4.00^\circ, &
SD_{\tilde{\theta}_\text{abs,max}} &= \sqrt{\frac{1.50}{4-1}} = 0.71^\circ \\
\bar{x}_\text{range} &= \frac{418.2+422.0+436.5+427.5}{4} = 426.1~\text{px}, &
SD_{x_\text{range}} &= \sqrt{\frac{189.33}{4-1}} = 7.9~\text{px}
\end{align*}

\noindent{\large\bfseries Calculations for disturbances around stable point 2:}\par
\begin{align*}
\bar{t}_{d,1} &= \frac{115.182+115.182+115.215+115.182}{4} = 115.190~\text{s}, &
SD_{t_d,1} &= \sqrt{\frac{0.000817}{4-1}} = 0.017~\text{s} \\
\bar{t}_{d,2} &= \frac{116.049+116.016+116.049+116.049}{4} = 116.041~\text{s}, &
SD_{t_d,2} &= \sqrt{\frac{0.000817}{4-1}} = 0.017~\text{s} \\
\bar{t}_{d,3} &= \frac{116.850+116.917+116.883+117.883}{4} = 117.133~\text{s}, &
SD_{t_d,3} &= \sqrt{\frac{0.7517}{4-1}} = 0.501~\text{s} \\
\bar{t}_{d,4} &= \frac{117.784+117.784+117.851+117.784}{4} = 117.801~\text{s}, &
SD_{t_d,4} &= \sqrt{\frac{0.003367}{4-1}} = 0.034~\text{s}
\end{align*}

\begin{align*}
\bar{\tilde{\theta}}_{\text{max},1} &= \frac{3.2+3.5+3.3+4.5}{4} = 3.63^\circ, &
SD_{\theta,1} &= \sqrt{\frac{1.0675}{4-1}} = 0.60^\circ \\
\bar{\tilde{\theta}}_{\text{max},2} &= \frac{6.7+4.1+6.4+6.8}{4} = 6.00^\circ, &
SD_{\theta,2} &= \sqrt{\frac{4.9000}{4-1}} = 1.28^\circ \\
\bar{\tilde{\theta}}_{\text{max},3} &= \frac{5.9+6.2+6.5+5.5}{4} = 6.03^\circ, &
SD_{\theta,3} &= \sqrt{\frac{0.5475}{4-1}} = 0.43^\circ \\
\bar{\tilde{\theta}}_{\text{max},4} &= \frac{5.5+4.4+4.9+4.3}{4} = 4.78^\circ, &
SD_{\theta,4} &= \sqrt{\frac{0.9075}{4-1}} = 0.55^\circ
\end{align*}

\begin{align*}
\Delta\bar{x}_{\text{max},1} &= \frac{75.3+94.2+89.8+103.1}{4} = 90.6~\text{px}, &
SD_{x,1} &= \sqrt{\frac{403.94}{4-1}} = 11.6~\text{px} \\
\Delta\bar{x}_{\text{max},2} &= \frac{92.9+89.3+99.1+99.6}{4} = 95.2~\text{px}, &
SD_{x,2} &= \sqrt{\frac{74.67}{4-1}} = 5.0~\text{px} \\
\Delta\bar{x}_{\text{max},3} &= \frac{104.3+98.2+99.9+104.5}{4} = 101.7~\text{px}, &
SD_{x,3} &= \sqrt{\frac{30.09}{4-1}} = 3.2~\text{px} \\
\Delta\bar{x}_{\text{max},4} &= \frac{59.6+60.9+94.4+64.1}{4} = 69.8~\text{px}, &
SD_{x,4} &= \sqrt{\frac{820.89}{4-1}} = 16.5~\text{px}
\end{align*}

\begin{align*}
&\textbf{1.4:} & t_r &= 118.785-117.784 = 1.001~\text{s}, &
&\textbf{2.4:} & t_r &= 118.819-117.784 = 1.035~\text{s} \\
&\textbf{3.4:} & t_r &= 118.785-117.851 = 0.934~\text{s}, &
&\textbf{4.4:} & t_r &= 118.785-117.784 = 1.001~\text{s}
\end{align*}

\begin{align*}
\bar{t}_r &= \frac{1.001+1.035+0.934+1.001}{4} = 0.993~\text{s}, &
SD_{t_r} &= \sqrt{\frac{0.005373}{4-1}} = 0.042~\text{s}
\end{align*}

\end{document}